\chapter{相关技术与理论基础}

\section{MRI 医学图像与数据特点}
MRI 通过核磁共振信号获取人体组织信息。与 CT 相比，MRI 对软组织的显示更细，也没有电离辐射，因此常用于脑部、脊柱和关节等部位的观察。本文使用脑部 MRI 数据，主要关注切片灰度、脑部轮廓和层间连续关系。

原始 MRI 数据通常不是一张普通图片，而是一个三维体数据。把体数据沿某个方向切开后，才能得到一张张二维切片。单张切片只能显示某一层的结构，连续切片才能体现脑部结构从上到下的变化。因此，本平台不能只生成孤立图片，还需要把切片顺序、空间位置和后续三维展示一起考虑。

MRI 与自然图像的差别比较明显。自然图像通常有丰富颜色和纹理，而 MRI 多为灰度图。它的边界、亮暗层次和组织轮廓更重要。不同设备、扫描参数和预处理方式也会影响灰度分布。评价这类生成结果时，不能只看“像不像”，还要说明结构是否合理、灰度是否稳定，以及结果可以用于什么场景。

\section{NIfTI 数据格式与预处理方法}
NIfTI 是神经影像领域常用的数据格式，文件后缀通常为 \texttt{.nii} 或 \texttt{.nii.gz}。它不仅保存体素灰度值，还可以保存体素尺寸、空间方向和仿射矩阵等信息\cite{nifti2004}。本文使用的 IXI 数据集以 NIfTI 格式提供脑部 MRI 数据，便于从体数据中提取切片，也便于后续组织连续序列\cite{ixiDataset}。

本项目的预处理目标比较直接：把原始体数据整理成 StyleGAN2 可以训练的 PNG 图片。系统先读取 NIfTI 文件，再沿轴向取切片。头顶部、颈部和边缘位置通常包含较多背景，因此脚本会优先保留体数据中间部分。灰度归一化时，系统使用 2\% 到 98\% 百分位截断，再把像素值映射到 $[0,255]$。这样可以减少少量异常亮点对整体灰度范围的影响。

完成灰度处理后，系统统一切片分辨率，并过滤均值、标准差和非零像素比例过低的图像。最终保存下来的 PNG 切片既用于 StyleGAN2 训练，也可以作为后续可视化和三维序列组织的基础。这个过程看起来只是数据整理，但它直接影响模型能否学到稳定的脑部轮廓。

\section{生成对抗网络基础}
GAN 由生成器 $G$ 和判别器 $D$ 组成。生成器从随机潜变量 $z$ 出发生成样本，判别器判断输入样本是真实数据还是生成数据。两者通过对抗训练共同优化。经典目标函数可写为：

\begin{equation}
\min_G \max_D V(D,G)=\mathbb{E}_{x\sim p_{data}(x)}[\log D(x)] + \mathbb{E}_{z\sim p_z(z)}[\log(1-D(G(z)))] .
\end{equation}

在理想状态下，生成器可以逐步学习真实数据分布，判别器也难以区分真实样本和生成样本。GAN 的优势是能够生成较清晰的图像，但训练过程容易出现模式崩溃、判别器过拟合和样本多样性不足等问题。医学影像数据规模有限，普通 GAN 训练往往不够稳定。本课题采用 StyleGAN2-ADA，主要考虑到它在有限数据训练场景中已有较成熟的工程实现。

\section{StyleGAN2 与 ADA 技术}
StyleGAN 系列方法先把输入潜变量映射到中间潜空间，再通过风格调制控制不同层级的图像特征\cite{karras2019stylegan}。这种结构使生成过程更容易控制。StyleGAN2 在此基础上改进了归一化方式，减少了生成图像中的伪影，并使用路径长度正则化提升潜空间映射稳定性\cite{karras2020stylegan2}。

StyleGAN2-ADA 的核心是自适应判别器增强（Adaptive Discriminator Augmentation，ADA）。训练数据较少时，判别器可能快速记住训练样本，生成器获得的有效训练信号会随之减弱。ADA 根据训练过程中的统计信号调整增强强度，从而降低判别器过拟合风险\cite{karras2020ada}。本课题的 MRI 切片来自有限数量的体数据。切片数量虽然较多，但病例来源仍有限，ADA 因而有助于提升训练稳定性。

\section{3D Gaussian Splatting 技术}
3D Gaussian Splatting 是一种三维场景表示和渲染方法。它用许多三维高斯基元表示空间中的结构。每个高斯都带有位置、尺度、方向、不透明度和颜色等参数。与体素网格相比，高斯点可以分布在连续空间中；与一些隐式表示相比，它更容易把结果导出并在前端进行预览\cite{kerbl20233dgs}。

对于第 $i$ 个三维高斯，可将其中心记为 $\mu_i$，协方差矩阵记为 $\Sigma_i$，颜色和不透明度分别记为 $c_i$ 与 $\alpha_i$。其空间影响范围可表示为：

\begin{equation}
G_i(x)=\exp\left(-\frac{1}{2}(x-\mu_i)^{T}\Sigma_i^{-1}(x-\mu_i)\right).
\end{equation}

渲染时，系统把三维高斯投影到目标视角，再根据深度、不透明度和颜色合成图像。训练时，模型比较渲染图像和目标图像，并根据差异更新高斯参数。原始 3DGS 通常依赖多视角照片和相机位姿。MRI 切片没有这种真实相机拍摄过程，它更接近按固定间距排列的灰度图层。因此，本文需要自己组织点云、切片序列和规则相机参数。

在本项目中，3DGS 的定位是三维表达和平台展示。系统先得到连续 MRI 切片，再按照 $z$ 轴顺序堆叠。非背景像素会被转换为空间点，灰度值则写入颜色字段，最终形成 PLY 初始点云。这个点云一方面作为 3DGS 训练起点，另一方面也能在前端直接展示，让用户看到二维切片怎样过渡到三维结构。

\section{FastAPI、React 与前后端分离技术}
本平台采用 FastAPI + React 的前后端分离架构。FastAPI 负责把预处理、生成、任务状态、图像资产和三维资产封装成接口。React 负责页面交互和结果展示。这样划分后，训练脚本不需要直接暴露给用户，前端也不需要了解底层命令行参数。

后端使用 Python 实现，主要调用 PyTorch、PIL、NumPy 和 nibabel 等库。耗时较长的训练和点云处理任务由后端启动子进程执行，任务日志和输出路径再返回给前端。对于需要实时查看状态的部分，后端通过 WebSocket 推送任务和日志变化。这样可以避免前端页面长时间卡住。

前端使用 TypeScript、React 和 Vite 构建。仪表盘、生成工作台、医学影像查看器、任务中心、三维展示和 AI 助手都作为页面或组件组织。前端通过统一 API 客户端访问后端接口，并使用 Cornerstone3D、Three.js、vtk.js 和 ECharts 完成切片浏览、点云展示、体渲染和指标图绘制。对用户来说，平台呈现的是一套网页操作流程；对系统内部来说，计算、接口和页面展示仍保持清楚分工。
